\documentclass[11pt]{letter}
\usepackage[utf8]{inputenc}
\usepackage[T1]{fontenc}
\usepackage{hyperref}
\usepackage[margin=1in]{geometry}
\signature{Zhongchang Huang\\Independent Researcher, Nanjing, China}
\date{June 5, 2026}

\begin{document}

\begin{letter}{Editor\\Physical Review B\\American Physical Society}

\opening{Dear Editor,}

We submit for your consideration our manuscript, ``Two-Exponent Characterization of Nonequilibrium Steady States: Coherence Decay and Transport in Power-Law Hopping Chains with Bulk Dephasing,'' for publication as a Regular Article in Physical Review B.

Transport in boundary-driven power-law hopping chains is typically characterized by a single exponent $\mu$ governing current decay with system size. The off-diagonal sector of the steady-state correlation matrix carries quantum coherence between distinct sites, but its spatial decay structure has not been characterized as a scaling phenomenon. We provide a self-consistent two-exponent characterization of the nonequilibrium steady state (NESS), computing both a coherence decay exponent $\beta(\alpha,\gamma_\phi)$ from midchain off-diagonal correlations and the transport exponent $\mu(\alpha,\gamma_\phi)$ from the boundary current---both within the same numerical framework.

Three results anchor the manuscript:

\begin{enumerate}
\item $\beta$ exhibits a two-regime large-$L$ structure. Using exact Lindblad solutions for 125 parameter combinations up to $L=64$, extended with sparse GMRES to $L=128$ for four dephasing slices and $L=256$ for $\gamma_\phi=0.5$, we find that $\beta$ descends from a long-range boundary-dominated sector at small $\alpha$ and crosses into a $\gamma_\phi$-dependent near-diffusive plateau at larger $\alpha$. A boundary-layer analysis quantifies the convergence: $\delta \sim L^{1.5}$ at $\alpha=1.1$ (boundary-dominated even at $L=256$), versus $\delta \sim L^{-0.57}$ at $\alpha=1.7$ (bulk-diffusive). $L=256$ data confirm the plateau is not a finite-size artifact: $|\beta_{128\to256} - \beta_{64\to128}| < 0.02$ for $\alpha \ge 1.5$.

\item $\beta$ and $\mu$ respond differently to parameter variations, confirming they probe distinct Liouvillian sectors. We compute $\mu$ self-consistently from the boundary current $J = \Gamma_L(f_L-D_1)$, finding $\mu = 0.285$--$0.901$ at $\gamma_\phi=0.5$---systematically below the dephasing-free formula $\mu=2\alpha-2$, as expected when bulk dephasing suppresses transport. Both $\beta$ and $\mu$ increase with $\gamma_\phi$ at fixed $\alpha$, but with different functional sensitivities ($\beta$ by $\sim$54\%, $\mu$ by $\sim$87\% from $\gamma_\phi=0.1$ to 2.0 at $\alpha=1.5$), demonstrating they are not locked in a fixed algebraic relation. The $(\mu,\beta)$ plane resolves transport-coherence regimes invisible to either exponent alone.

\item $\beta$ is robust under boundary-coupling variations. A factor-of-four change in $\Gamma$ shifts $\beta$ by under 7\%, while the exact $\Delta f$-independence follows from the linearity of the Lyapunov equation. The weak $\Gamma$-dependence is interpreted as a boundary-fixed-line effect.
\end{enumerate}

We prove, via the homogeneous Lyapunov system argument, that off-diagonal correlations are purely imaginary for arbitrary power-law hopping---generalizing the nearest-neighbor result of Bhat and \v{Z}nidari\v{c} [Phys.\ Rev.\ B~111, 174306 (2025)]. The proof is structural and independent of the $\beta$-exponent analysis.

We have been careful to distinguish what is derived from what is calibrated. The $1/\alpha$ leading scaling coordinate follows from the fractional diffusion kernel and fixed-point constraints; amplitudes are calibrated numerically. Full model selection (AIC$_c$ for all five $\gamma_\phi$ values) is reported in the Supplemental Material. A first-order commutator expansion shows $\beta \to \nu$ (the occupation-gradient exponent) in the thermodynamic limit; the finite-size correction $\kappa \neq 1$ distinguishes $\beta$ from $\nu$ at experimentally accessible $L$, and the two-regime large-$L$ structure contains physics beyond the thermodynamic limit. All raw data, analysis scripts, and figure-generation code are archived at Zenodo (DOI: \url{https://doi.org/10.5281/zenodo.20553343}).

The manuscript has not been published elsewhere and is not under consideration by another journal. We have no conflicts of interest to declare.

We believe this self-consistent two-exponent characterization---combining exact numerics up to $L=256$, a boundary-layer scaling analysis, and direct experimental relevance to trapped-ion and Rydberg platforms---will interest the broad readership of Physical Review B.

\closing{Sincerely,}

\end{letter}
\end{document}
